\documentclass[12pt,a4paper]{article}
\usepackage{geometry}
\geometry{left=2.5cm,right=2.5cm,top=2.0cm,bottom=2.5cm}
\usepackage[english]{babel}
\usepackage{amsmath,amsthm}
\usepackage{amsfonts}
\usepackage[longend,ruled,linesnumbered]{algorithm2e}
\usepackage{fancyhdr}
\usepackage{ctex}
\usepackage{array}
\usepackage{listings}
\usepackage{color}
\usepackage{graphicx}

\begin{document}


\title{
{《优化理论与算法》第13次作业
}
}
\date{}

% \author{
% 姓名：\underline{你的名字}~~~~~~
% 学号：\underline{你的学号}~~~~~~}

\maketitle
\begin{itemize}
	\item 作业截止于{\bf 周五（6/12/2025）}，在Canvas系统中提交PDF或照片文件
	\item 作业题目的tex文件可以在Canvas中下载
	\item 鼓励相互讨论，但不要抄袭.
\end{itemize}
%%%%%%%%%%%%%%%%%%%%%%%%%%%%%%%%%%%%%%%%%%%%%%%%%%%%%%%%%%%%%%%%%
%%%%%%%%%%%%%%%%%%%%%%%%%%%%%%%%%%%%%%%%%%%%%%%%%%%%%%%%%%%%%%%%%

\section{随机梯度下降法}

\paragraph{1.1} 比较所学随机优化方法各自的优势

\paragraph{1.2} 常见的学习率调节方法有哪些

\paragraph{1.3}
\textbf{问题构建:} 在二维平面中，考虑以下设定：
\begin{itemize}
    \item 已知锚点位置：$\mathbf{a}_i \in \mathbb{R}^2,\ i=1,\dots,m$；
    \item 未知节点位置：$\mathbf{x}_j \in \mathbb{R}^2,\ j=1,\dots,n$（待估计）。
\end{itemize}

我们可以获得带噪声的\textbf{平方距离}测量值 $d_{ij}$，表示锚点 $i$ 到未知节点 $j$ 的平方距离测量（而非直接距离）。定义如下非线性最小二乘目标函数：

\[
\min_{\mathbf{x}_1,\dots,\mathbf{x}_n} \; F(\mathbf{x}) = \sum_{i=1}^{m}\sum_{j=1}^{n} \bigl( \|\mathbf{a}_i - \mathbf{x}_j\|_2^2 - d_{ij} \bigr)^2 .
\]
来求解未知节点的位置。

\textbf{数据生成:} 
\begin{enumerate}
    \item 参数设置
\begin{itemize}
    \item 锚点数量：$m = 100$；
    \item 未知节点数量：$n = 50$；
    \item 坐标范围：$[0, 10] \times [0, 10]$；
    \item 噪声标准差：$\sigma = 0.1$；
    \item 随机种子固定（例如 $42$）以保证结果可重复。
\end{itemize}

\item 生成真实位置
\begin{itemize}
    \item 锚点坐标 $\mathbf{a}_i$：在 $[0,10]^2$ 内均匀随机生成。
    \item 未知节点真实坐标 $\mathbf{x}_j^{\text{true}}$：在 $[0,10]^2$ 内均匀随机生成。
\end{itemize}

\item 计算带噪声的平方距离测量
对于每一对 $(i,j)$，$i=1,\dots,m$，$j=1,\dots,n$：
\begin{itemize}
    \item 真实平方距离：$d_{ij}^{\text{true}} = \|\mathbf{a}_i - \mathbf{x}_j^{\text{true}}\|^2$；
    \item 测量值：$d_{ij} = d_{ij}^{\text{true}} + \epsilon_{ij}$，其中 $\epsilon_{ij} \sim \mathcal{N}(0, \sigma^2)$ 且相互独立。
\end{itemize}
测量值 $d_{ij}$ 可能为负值（由于噪声），这在实际数据中是允许的，仅作为优化算法的输入。

\end{enumerate}

\textbf{输出要求：}  比较随机梯度下降法、随机梯度下降法+动量、随机梯度下降法+方差缩减，Adam优化方法在该问题的求解效果。包含代码实现，绘制相关图像，总结观察等。



\section{内点法}

\paragraph{2.1}
线性规划（标准型）的原始对偶中心路径是什么? 

\paragraph{2.2}
考虑如下线性规划
\begin{align*}
\max & \,\,\,\, 2x\\
\text{s.t. } & x \leq 3 \\
& x \leq 5 \\
& x \geq 0
\end{align*}


\begin{enumerate}
	\item[(a)] 给定障碍参数$\mu=1$，写出原始对偶中心路径的形式。
	
	\item[(b)] 考虑原始对偶可行解 $x=2$ 和 $y=(1,2)$。求a）问对应的牛顿步。
	
	\item [(c)] 对比初始解$x=2, y=(1,2)$的对偶间隙 与一次牛顿迭代后的对偶间隙。
	
\end{enumerate}


\paragraph{2.3}
比较单纯形法与内点法各自的优势


\end{document}